% GENERATED, edit the French unit or the translation table
\mxcorrections{6}\label{C06-CORR}
\textbf{Teacher use.} The answers below concern the tasks not worked in the lesson. First have students state their method; give graded help before the solution. The durations are estimates to be calibrated in class.
\subsection*{Diagnostic and remediation}
\begin{enumerate}
\item $6$. \textit{Action:} Revisit the rate $2a+h$.
\item $\R\setminus\{2\}$. \textit{Action:} Set the denominator to zero.
\item Positive outside $\intff{-2}{1}$, negative between the roots, zero at the endpoints. \textit{Action:} Build two sign lines.
\item $x^3-2x^2+x-2$. \textit{Action:} Distribute each term without omission.
\item They are equal. \textit{Action:} Substitute -x in brackets.
\item $y=2$. \textit{Action:} Revisit the tangent equation and simplify.
\end{enumerate}
\subsection*{Investigation activity}
$x^3$ is strictly increasing with a horizontal tangent at 0. For a differentiable increasing function, the rates on either side are positive or zero, so the derivative cannot be negative.
\subsection*{Immediate applications}
\begin{enumerate}
\item $5x^4$ on $\R$; $-\dfrac{2}{x^3}$ for $x\ne0$; $\dfrac{1}{2\sqrt x}$ for $x>0$.
\item $2x(3x-2)+3(x^2+1)=9x^2-4x+3$; $-\dfrac{2x}{(x^2+1)^2}$, on $\R$.
\item $f\prime=2x-4$, so decreasing up to 2 then increasing. Minimum $-3$ at 2; values at the endpoints 1 and $-2$, so maximum 1 at 0.
\item First even, second odd; difference $x^2+1-2x=(x-1)^2\geq0$, equality at 1.
\item $A(x)=x(6-x)$ on $\intff{0}{6}$; $A\prime=6-2x$. Maximum 9 m$^2$ at 3.
\end{enumerate}
\subsection*{Exercises: method and results}
\paragraph{Exercise 6.1 -- A collection of derivatives.} $12x^3-2$; $2x(x-1)+(x^2+2)=3x^2-2x+2$; $\dfrac{[2(x+3)-(2x+1)]}{(x+3)^2}=\dfrac{5}{(x+3)^2}$ for $x\ne-3$.
\paragraph{Exercise 6.2 -- A complete table.} $f\prime=3(x-1)(x-3)$; signs $+,-,+$. Values $f(0)=1$, $f(1)=5$, $f(3)=1$, $f(4)=5$. Minimum 1 at 0 and 3; maximum 5 at 1 and 4. The endpoints count.
\paragraph{Exercise 6.3 -- A quotient.} $f\prime =1-\dfrac{4}{x^2}=\dfrac{(x-2)(x+2)}{x^2}$. Sign negative before 2, positive after. Minimum 4 at 2; for $x>0$, $x+\dfrac{4}{x}\geq 4$, equality only at 2.
\paragraph{Exercise 6.4 -- Comparing two curves.} $f-g=x(x-1)(x+1)$. Signs: negative on $\intoo{-\infty}{-1}$, positive on $\intoo{-1}{0}$, negative on $\intoo{0}{1}$, positive on $\intoo{1}{+\infty}$. Intersection at $-1,0,1$. We can make the table of the three factors without differentiating.
\paragraph{Exercise 6.5 -- A box.} $V=x(12-2x)^2=4x(6-x)^2$. $V\prime=12(x-2)(x-6)$, positive on $\intfo{0}{2}$, negative on $\intoo{2}{6}$. $V(0)=V(6)=0$, $V(2)=128$ cm$^3$. Maximum for x=2 cm; box $8\times8\times2$.
\paragraph{Exercise 6.6 -- Stationary without an extremum.} $f\prime=3x^2$, so value 0 at 0. In every neighbourhood of 0, negative x give $f(x)<0$ and positive x give $f(x)>0$; no extremum.
\paragraph{Exercise 6.7 -- Comparing strategies.} Square: $(x-3)^2+3\geq3$. Derivative: $2x-6$, minimum at 3, value 3. The first avoids a complete study of variation; the second illustrates the link between chapters.
\paragraph{Exercise 6.8 -- Delayed review.} It decreases on $\intoo{-\infty}{0}$ and $\intoo{0}{+\infty}$. It does not decrease on their union: $-1<1$ but $f(-1)=-1<1=f(1)$. The domain is not an interval.
\paragraph{CQFD, exercise 12 p.~156, questions 1 and 2.} $f(x)-g(x)=\dfrac12x^2-1=\dfrac12(x-\sqrt2)(x+\sqrt2)$. The curve of f is strictly above that of g for $x<-\sqrt2$ or $x>\sqrt2$, below between these two numbers; equality at the x-coordinates $\pm\sqrt2$.
\paragraph{CQFD, exercise 16 p.~156, questions f and g.} For $f(x)=x^3$, we can choose $\intoo{-\infty}{0}$: $f(x)<0$ and the function is strictly increasing. For $g(x)=x^2-4$, we can choose $\intfo{0}{2}$: $g(x)<0$ and the function is strictly increasing on this interval. Subintervals that are not reduced to a point also work.
\subsection*{MCQ: answers and reasons}
\textbf{MCQ 1.}
\begin{enumerate}
\item B. Integer power.
\item A. Rate of the reciprocal.
\item D. Product rule.
\item C. Sign of the derivative and variation.
\item B. Nonzero denominator.
\end{enumerate}
\textbf{MCQ 2.}
\begin{enumerate}
\item C. $f(-x)=-f(x)$.
\item A. The other conclusions require more.
\item D. Formula and domain.
\item B. Reached at 1.
\item C. Derivative $2x_0=0$.
\end{enumerate}
\subsection*{True or false}
\begin{enumerate}
\item False: $f(x)=\dfrac{1}{x}>0$ on $\intoo{0}{+\infty}$, but $f\prime(x)<0$.
\item True, with the hypotheses stated; at the endpoints, the conclusion is not imposed.
\item False without sign hypotheses: on $\intff{-2}{-1}$, $u(x)=v(x)=x$ are increasing but $x^2$ is decreasing.
\end{enumerate}
\paragraph{Differentiated teaching.} For a technical block, offer only the first step (domain, factor or reference sample space), then withdraw the help. For an independent student, ask for a second method and the treatment of a limiting case. Orally, allow an initial reformulation in the language the student has mastered, then require a precise mathematical sentence in French.
